\chapter*{Introduction Générale}
\addcontentsline{toc}{chapter}{Introduction Générale}
\markboth{Introduction Générale}{}

L'essor fulgurant des technologies de communication numérique a profondément transformé les modes d'organisation et de participation aux événements professionnels et académiques. Les plateformes de visioconférence sont désormais au c\oe{}ur de la vie économique et sociale, rassemblant des milliers de participants à travers le monde en temps réel. Cette évolution s'accompagne cependant de risques croissants en matière de sécurité informatique : interception des flux de données, usurpation d'identité, violation de la vie privée et non-conformité aux réglementations en vigueur constituent autant de menaces auxquelles les opérateurs de ces plateformes doivent répondre.

C'est dans ce contexte que s'inscrit le présent projet de fin d'études, intitulé \textit{« Conception et Implémentation d'une Infrastructure de Sécurité pour une Plateforme Événementielle en Ligne Optimisée par l'IA »}. L'ambition de ce travail est double : d'une part, bâtir une infrastructure de sécurité multicouche, conforme aux meilleures pratiques de la cybersécurité et aux exigences du Règlement Général sur la Protection des Données (RGPD) ; d'autre part, enrichir l'expérience des participants grâce à des services d'intelligence artificielle, notamment la transcription et la traduction automatiques en temps réel.

La plateforme proposée répond à un besoin métier concret : permettre à un organisateur de planifier, configurer et animer des événements en ligne, en s'appuyant sur des templates personnalisés et une synchronisation avec le progiciel de gestion Odoo. Les participants bénéficient quant à eux d'une interface accessible, sécurisée et multilingue, renforçant ainsi l'inclusivité et la portée internationale des événements hébergés.

Sur le plan méthodologique, ce projet suit un cycle de développement structuré, articulé autour des phases classiques du génie logiciel : analyse des besoins, conception architecturale, implémentation et validation. Ce rapport rend compte de l'intégralité de cette démarche, en présentant les choix techniques adoptés, les mécanismes de sécurité mis en \oe{}uvre et les résultats obtenus.

Le présent document est organisé en cinq chapitres principaux :

\begin{itemize}
    \item \textbf{Chapitre 1 — Cadre du Projet} : ce chapitre présente l'organisme d'accueil (TEK-UP), le contexte et la problématique du projet, ses objectifs et la méthodologie agile SCRUM adoptée.
    \item \textbf{Chapitre 2 — État de l'Art} : ce chapitre dresse un panorama des domaines scientifiques et technologiques mobilisés — cybersécurité, gestion des incidents, supervision réseau, IA, SOC et SIEM.
    \item \textbf{Chapitre 3 — Spécification des Besoins} : ce chapitre identifie les acteurs, formalise les besoins fonctionnels et non fonctionnels, présente les études comparatives et justifie les choix technologiques.
    \item \textbf{Chapitre 4 — Conception de l'Architecture} : ce chapitre présente l'architecture globale du système, la modélisation UML et les mécanismes de sécurité (STRIDE, RBAC, RGPD by Design).
    \item \textbf{Chapitre 5 — Implémentation et Déploiement} : ce chapitre décrit la mise en \oe{}uvre technique, les configurations de sécurité appliquées et les résultats des tests sur infrastructure OVHcloud.
\end{itemize}
